% Preamble additions for the Architecture Vision document, injected by pandoc
% with -H. Deliberately NOT ../../../../preamble.tex: that one is the NULP lab
% preamble — it sets polyglossia to Ukrainian, names captions «Рис.» and
% «Таблиця», loads minted for code listings this document has none of, and ends
% by \input-ing titlepage.tex. This document is English, has no NULP title page,
% and reproduces the Word template's own cover from \maketitle, redefined below.
%
% The whole file exists to make the PDF match "Architecture Vision
% Template_masters.docx" page for page: Letter paper, the template's margins,
% running head and foot, heading sizes, table style and list indents are all
% taken from that document's styles.xml rather than chosen. Check a change
% against the original before making it.
%
% BUILD: LuaLaTeX, via scripts/render-md-report.sh. Required by fontspec below.

% Open Sans at 11pt is the Word template's Normal style, and it carries the
% Cyrillic the author block needs. Both families below ship with TeX Live and
% are found by filename through kpathsea, so no Path= and no fontconfig entry —
% naming them the way fontconfig would ("Open Sans") is what fails here.
\setmainfont{OpenSans-Regular.ttf}[
  BoldFont=OpenSans-Bold.ttf,
  ItalicFont=OpenSans-Italic.ttf,
  BoldItalicFont=OpenSans-BoldItalic.ttf,
]
\setmonofont{DejaVu Sans Mono}[Scale=0.82]

% Headings are set in the body face at regular weight — deliberately, with no
% separate heading family. The template's Heading 1-4 styles name "Proxima Nova
% Bl", which is commercial and installed nowhere in this toolchain or on the
% machines the document is read on, so Word and every converter substitute the
% body face: the template as anyone actually sees it has light headings whose
% only weight contrast is the bold section number, which the styles apply to
% the list number alone and not to the heading text. A black-weight stand-in
% reproduces the style sheet's intent but not the document, and the document is
% what has to match.
% US Letter with 1in margins all round, matching the template's w:pgMar
% (1440 twips on every side). The paper size also comes from the class option
% passed by the build script; geometry alone would not correct it. An earlier
% revision used A4, which fitted ~16% more text per page and so paginated
% differently from the Word original — the document is a filled-in template and
% has to lay out like one.
\usepackage[letterpaper,left=1in,right=1in,top=1in,bottom=1in,
            headsep=14pt,footskip=30pt]{geometry}
\usepackage{setspace}
% A touch looser than the template's single-spaced Normal style: the prose here
% is denser than the issued template's placeholder text, and at 11pt on a
% 6.5in measure the extra leading is what keeps a full page readable. Paragraph
% separation stays the 6pt gap set further down, not the leading, so widening
% this changes line rhythm only and not the gaps between blocks.
\setstretch{1.15}

\usepackage{microtype}

% Loaded unconditionally: pandoc's template pulls graphicx in only when the
% document actually contains an image, but the build script always appends a
% \graphicspath, and that is undefined without it. Re-loading is a no-op.
\usepackage{graphicx}

% Figures sit exactly where the Markdown puts them, as the Word original's
% images do. Without this LaTeX floats them to the next convenient page, which
% in a document this sparse means a diagram drifting a page away from the
% section that introduces it, with a half-empty page left behind.
% Rewriting the environment to pass [H] explicitly, rather than setting
% \fps@figure: float.sty only honours H when it is given as the optional
% argument, and pandoc emits \begin{figure} without one.
\usepackage{float}
\let\origfigure\figure
\let\endorigfigure\endfigure
\renewenvironment{figure}[1][]{\origfigure[H]}{\endorigfigure}

% The six-part quality attribute scenario in 3.4 is the one figure drawn rather
% than imported: in the Word template it is a drawing canvas whose shapes
% autofit their text, and every converter clips it ("Respons", "credentials"
% cut mid-word). Redrawing it keeps the template's layout with text that fits.
\usepackage{tikz}
\usetikzlibrary{positioning,arrows.meta,fit}

% colortbl backs the black header row that scripts/md-report-tables.lua emits;
% longtable and array supply the environment and the p-column modifiers it uses.
\usepackage{etoolbox}
\usepackage{longtable}
\usepackage{array}
\usepackage{colortbl}

% The lab preamble's labelindent=1.27cm leaks into table cells and pushes
% bullets against the cell edge; every table in this template has list cells.
% 18pt is the Word list indent (0.25in to the bullet, 0.5in to the text), and
% the 6pt item separation is the Normal style's paragraph spacing, which Word
% applies between list items too.
\usepackage{enumitem}
\setlist{topsep=6pt,itemsep=6pt,parsep=0pt,partopsep=0pt}
\setlist[itemize]{label=\textbullet,labelindent=18pt,leftmargin=*,labelsep=7pt}
\setlist[enumerate]{labelindent=18pt,leftmargin=*,labelsep=7pt}
% Inside a table cell the indent has nowhere to go, so lists reset to flush.
% \parskip goes with them: a cell opens a paragraph before the list starts, so
% the 6pt body spacing would land on top of the list's own and push the first
% bullet off the top of the cell, out of line with its neighbours.
% Keyed by type, not the plain \setlist: enumitem merges the two and the
% type-specific declaration above wins every key it names, so a generic reset
% here would leave the 18pt label indent in place inside the cells.
\AtBeginEnvironment{longtable}{%
  \setlength{\parskip}{0pt}%
  \setlist[itemize]{topsep=0pt,partopsep=0pt,itemsep=4pt,parsep=0pt,
    label=\textbullet,labelindent=0pt,leftmargin=*,labelsep=5pt}%
  \setlist[enumerate]{topsep=0pt,partopsep=0pt,itemsep=4pt,parsep=0pt,
    labelindent=0pt,leftmargin=*,labelsep=5pt}}

% labelsep=period gives "Figure 3.1. Caption" rather than "Figure 3.1: Caption".
% The template's Caption style: centred, italic, 10pt.
\usepackage[labelsep=period]{caption}
\captionsetup{font={it,small},justification=centering}
\captionsetup[table]{justification=raggedright,singlelinecheck=false}

% Headings are numbered four levels deep: the template numbers the view pattern
% too, as 3.2.1.1 View Context / 4.1.1.1 Intent. --top-level-division=section
% maps '#' to \section, so the figure and table counters must be re-tied to
% section — left on chapter they would print "0.1".
\setcounter{secnumdepth}{4}
% extreport still counts chapters, so an unused chapter 0 would prefix every
% number: "0.1 Introduction". Re-root the section numbers at \section.
\renewcommand{\thesection}{\arabic{section}}
\renewcommand{\thesubsection}{\arabic{section}.\arabic{subsection}}
\renewcommand{\thesubsubsection}{\arabic{section}.\arabic{subsection}.\arabic{subsubsection}}
\renewcommand{\theparagraph}{\thesubsubsection.\arabic{paragraph}}
\numberwithin{figure}{section}
\numberwithin{table}{section}
\numberwithin{equation}{section}

% Contents page, matching the template's: every level dot-led, level 1 bold with
% a bold page number, the three levels stepping in from the left margin.
% Patched directly rather than through tocloft, which bails out with
% "\@starttoc has already been redefined" -- pandoc's template gets there first,
% and header-includes are injected at the end of the preamble, so tocloft can
% never win that race from here. report.cls already dot-leads \l@section
% because the top-level division is \section, so only the weights and the
% indents need setting.
\makeatletter
\renewcommand*\l@section[2]{\@dottedtocline{1}{0em}{1.6em}{\bfseries #1}{\bfseries #2}}
\renewcommand*\l@subsection{\@dottedtocline{2}{1.0em}{2.2em}}
\renewcommand*\l@subsubsection{\@dottedtocline{3}{2.4em}{3.0em}}
% \@dotsep is in math units; LaTeX's default 4.5 spaces the dots far wider than
% the Word leader, which is near-continuous.
\renewcommand*\@dotsep{1.3}
\makeatother

\usepackage{titlesec}
% Each top-level heading opens a page, as the Word template has it. titlesec
% runs \sectionbreak ahead of every \section, starred ones included, so this
% covers the unnumbered Revision History too; \clearpage rather than \newpage so
% a float cannot drift past the heading it belongs to.
\newcommand{\sectionbreak}{\clearpage}
% Sizes are the template's Heading 1-4: 24, 20, 16 and 12pt. The number is the
% only part set bold, matching the template — see the font note at the top.
% \chapter survives in this document only as the heading \tableofcontents emits,
% and it carries no number, so it stays unbolded throughout like the heading
% text beside it.
\titleformat{\chapter}[block]{\fontsize{24}{28.8}\selectfont}{}{0pt}{}
\titlespacing{\chapter}{0pt}{0pt}{16pt}
\titleformat{\section}[block]{\fontsize{24}{28.8}\selectfont\raggedright}{\bfseries\thesection}{0.6em}{}
\titleformat{\subsection}[block]{\fontsize{20}{24}\selectfont\raggedright}{\bfseries\thesubsection}{0.6em}{}
\titleformat{\subsubsection}[block]{\fontsize{16}{19.2}\selectfont\raggedright}{\bfseries\thesubsubsection}{0.6em}{}
% Non-starred \titlespacing is required: the starred form titlesec applies by
% default kills the indent of the first paragraph after a heading.
\titlespacing{\section}{0pt}{20pt}{12pt}
\titlespacing{\subsection}{0pt}{14pt}{6pt}
\titlespacing{\subsubsection}{0pt}{12pt}{4pt}
% The template's view pattern (Intent / Context / Representation / Element
% Catalog) sits at '####' — the template's own Heading 4, 12pt. It gets its own
% line and its own number (3.2.1.1), as in the Word original; a run-in form
% reads better but paginates and numbers differently, and this document has to
% match the template it is filled into.
\titleformat{\paragraph}[block]{\fontsize{12}{14.4}\selectfont\raggedright}{\bfseries\theparagraph}{0.6em}{}
\titlespacing{\paragraph}{0pt}{10pt}{4pt}

% ---------------------------- Title page ----------------------------
% The front page's layout is fixed template furniture, so it is written out here
% rather than assembled by \maketitle from meta.yaml: the logo, the title at
% 32pt (Title page1), the project and version block at 16pt and 14pt (Title
% page2/3), and the date at 10pt grey (Title page date). The four values the
% template leaves as placeholders are set once below; the version also feeds the
% running head, so a new revision changes one line. The cover carries no author,
% as the template has none: the submitter's name belongs in the Revision History
% table. \subtitle is neutralised for the same reason: pandoc emits it into the
% preamble whenever meta.yaml defines one, and it must not reach \@title.
\newcommand{\subtitle}[1]{}
\newcommand{\avproject}{eDiscovery and Litigation Document Review Platform}
% The business case never names the firm, so the cover describes it instead of
% inventing a name.
\newcommand{\avclient}{International Law Firm}
\newcommand{\avversion}{0.2}
% mm/dd/yyyy, the format of the template's own placeholder.
\newcommand{\avdate}{10/09/2026}
\renewcommand{\maketitle}{%
  \begin{titlepage}%
    \thispagestyle{empty}%
    \vspace*{-12mm}%
    \noindent\includegraphics[width=52mm]{softserve-logo.png}%
    \vspace*{66mm}%
    \par\noindent{\fontsize{32}{38.4}\selectfont ARCHITECTURE VISION}%
    \vspace*{12mm}%
    \par\noindent{\fontsize{16}{19.2}\selectfont \avproject}%
    \vspace*{8mm}%
    \par\noindent{\fontsize{14}{16.8}\selectfont Version \avversion}%
    \vspace*{4mm}%
    \par\noindent{\fontsize{14}{16.8}\selectfont For}%
    \vspace*{4mm}%
    \par\noindent{\fontsize{14}{16.8}\selectfont \avclient}%
    \vfill
    \par\noindent{\fontsize{10}{12}\selectfont\color{black!55}\avdate}%
    \vspace*{4mm}%
  \end{titlepage}%
  % The titlepage environment resets the counter to 1, which would put Contents
  % on page 1; the template counts the cover.
  \setcounter{page}{2}%
}

% Running head and foot, as the template carries them on every page but the
% title: the document name and version ranged right at the top, the
% confidentiality marking and the page number at the foot. Both come from the
% Word template's header/footer parts and are reproduced rather than invented.
\usepackage{fancyhdr}
\pagestyle{fancy}
\fancyhf{}
\fancyhead[R]{\footnotesize\textbf{ARCHITECTURE VISION}\\[-2pt]\scriptsize Version \avversion}
\fancyfoot[L]{\footnotesize SoftServe Confidential}
\fancyfoot[R]{\footnotesize\textbf{\thepage}}
\renewcommand{\headrulewidth}{0pt}
\renewcommand{\footrulewidth}{0pt}
\setlength{\headheight}{24pt}
% \chapter and \section* (the contents page) switch to `plain` on their own, so
% it has to carry the same header and footer or those pages would lose them.
\fancypagestyle{plain}{%
  \fancyhf{}%
  \fancyhead[R]{\footnotesize\textbf{ARCHITECTURE VISION}\\[-2pt]\scriptsize Version \avversion}%
  \fancyfoot[L]{\footnotesize SoftServe Confidential}%
  \fancyfoot[R]{\footnotesize\textbf{\thepage}}%
  \renewcommand{\headrulewidth}{0pt}%
}

% Block paragraphs, as in the Word template: no first-line indent, so the gap
% between paragraphs is the only thing separating them and cannot be 0 — pandoc
% leaves parindent at 0pt and would otherwise run two paragraphs together. 6pt
% is the template's Normal style, w:spacing w:after="120" twips.
\setlength{\parindent}{0pt}
\setlength{\parskip}{6pt plus 2pt minus 1pt}
\renewcommand{\arraystretch}{1.3}

% Hyphenation off, as in the lab template.
\hyphenpenalty=10000
\exhyphenpenalty=10000

% Ragged right, which is what the Word template's Normal style does: Word only
% justifies when told to, and this one is not. Without it LaTeX justifies, and
% with hyphenation off above all the slack lands in the interword spaces --
% visibly looser lines than the original, and a different line count per
% paragraph, which is half of why the pagination drifted. \raggedright has to
% be deferred: pandoc's template and the geometry/setspace setup both touch the
% paragraph shape after header-includes is read.
\AtBeginDocument{\raggedright}

% Every link in this document is an internal cross-reference, and pandoc sets
% linkcolor=black for those -- so without this they are typeset exactly like the
% surrounding prose and nothing on the page says "Business Case" is clickable.
% Underlining rather than colouring keeps the body a single colour, as the Word
% template's is, and survives printing.
%
% The kernel's \underline, not ulem's \uline: ulem is not in this TeX Live
% installation, and the document builds everywhere without it. The cost is that
% \underline boxes its argument, so an underlined cross-reference cannot break
% across a line -- with the ragged right margin set above a long one wraps whole
% to the next line, which costs a little raggedness and never an overfull line.
%
% \hyperref is patched rather than hyperref's internal \Hy@colorlink, which
% would have caught every link kind at once: the internal hook brackets the link
% text with two separate macro calls, and \underline needs its argument as one
% balanced group, so wrapping it there makes TeX scan past the closing brace
% looking for an \egroup that only appears inside a macro body. Patching the
% user-level command keeps the text an explicit argument. The two-argument
% \hyperref{url}{category}{name}{text} form is left alone -- pandoc emits only
% the \hyperref[label]{text} one for the Markdown cross-references this document
% is built from, and \href and \url appear nowhere in it.
%
% The patch has to be deferred the same way \raggedright is, and for a stronger
% reason: pandoc loads hyperref (through bookmark) after header-includes, so
% \hyperref does not exist yet while this file is read.
\makeatletter
\def\MDR@hyperref@opt[#1]#2{\MDR@hyperref[#1]{\underline{#2}}}
\AtBeginDocument{%
  % \NewCommandCopy, not \let: \hyperref is robust, so what \let would copy is
  % only the protection wrapper, whose body lives in a second control sequence
  % that the \DeclareRobustCommand below overwrites as well -- the saved copy
  % would then reach the new definition and recurse until LuaTeX wedges.
  \NewCommandCopy\MDR@hyperref\hyperref
  \DeclareRobustCommand\hyperref{%
    \@ifnextchar[{\MDR@hyperref@opt}{\MDR@hyperref}}%
}
\makeatother
